\section{The canonical singular Ricci-flat metric}\label{sec:canonical}

\subsection{Identification by the Monge--Amp\`ere equation}
Let $\sigma$ denote the holomorphic symplectic form on $\Mn_{\rm reg}$,
normalized so that its restriction to $\M$ is the form supplied by
hyperk\"ahler reduction. Since $\pi$ is a symplectic resolution, this form
extends nondegenerately to $\MG$.

\begin{theorem}\label{thm:canonical}
The current $T$ of Theorem~\ref{thm:global-current} is the canonical
Ricci-flat current in the ample class $\Lambda'$. It is a smooth
hyperk\"ahler metric on $\Mn_{\rm reg}$, and its restriction to $\M$ is
$\gL$. Moreover,
\begin{equation}\label{eq:volume}
 \Vol(\M,\gL)=\frac{(16\pi^4\Vol(X,g_X))^N}{N!}.
\end{equation}
\end{theorem}
\begin{proof}
The $L^2$ metric is hyperk\"ahler \cite[Theorem~2.1]{Itoh88}. Consequently,
on $\M$ its volume form satisfies
\[
 \Om^{2N}=b\,(\sigma\wedge\overline\sigma)^N
\]
for a positive constant $b$, with the conventional sign absorbed in the
positive measure on the right. The left side is the Bedford--Taylor
Monge--Amp\`ere measure of a continuous-potential current. It gives no mass
to proper analytic subsets. The symplectic measure on the right also gives
no mass to the boundary: on the resolution it is a smooth volume form.
The identity therefore holds on all of $\Mn$, in the sense of measures.
The cohomology class fixes $b$ by integration.

The space $\Mn$ is projective, canonical and has trivial canonical class.
The existence and uniqueness theorem of Eyssidieux--Guedj--Zeriahi
\cite[Theorem~7.5 and Corollary~7.6]{EGZ09} applies to the ample real
algebraic class $\Lambda'$. Its unique bounded-potential solution is $T$.
The same theorem makes $T$ smooth on $\Mn_{\rm reg}$. The parallel
symplectic structure agrees there with that on the dense open set $\M$,
so this regular metric is hyperk\"ahler.

Bounded potentials have full cohomological mass; hence
\begin{equation}\label{eq:intersection-volume}
 \Vol(\M,\gL)=\frac1{(2N)!}\int_{\MG}\Lambda^{2N}.
\end{equation}
Let $q$ be the Beauville--Bogomolov form normalized by the Mukai isometry.
The vector associated to $\mu_v(\alpha)$ is, up to sign,
$(0,\alpha,c\cdot\alpha/r)\in v^\perp$. It has square $\alpha^2$.
The Mukai isometry \cite{OGr97,Yos01} thus gives
\[
 q(\Lambda)=16\pi^4\int_X\omega^2
           =32\pi^4\Vol(X,g_X).
\]
The Fujiki relation for K3$^{[N]}$ type is
\[
 \int_{\MG}\Lambda^{2N}=
 \frac{(2N)!}{2^NN!}\,q(\Lambda)^N.
\]
Substitution in \eqref{eq:intersection-volume} proves \eqref{eq:volume}.
\end{proof}

The right side of \eqref{eq:volume} is also the volume of
$\Sym^N(X,4\pi^2g_X)$. This equality is an intersection-theoretic identity;
it does not assert an isometry with that symmetric product.

\subsection{The Quillen metric}
The preceding argument also controls a natural determinant metric. For a
large positive integer $k$, put
\begin{equation}\label{eq:quillen-line}
 \mathcal L=\big[\det Rp_*(R\mathcal Hom(\mathbb E,\mathbb E)\otimes
                  (H^k-H^{-k}))\big]^{-1}.
\end{equation}
This expression is read on the bundle locus for its Quillen metric and
in perfect K-theory for its global determinant class.
A fixed trivial-rank term may be included in this virtual bundle without
affecting its varying curvature. Let $m$ be the similitude of the chosen quasi-universal sheaf. The
actual determinant line retains the factor $m^2$ in its curvature and
class; positive powers clear rational descent data, not this factor.

\begin{corollary}\label{cor:quillen}
On $\M$, a positive power of the Quillen metric on $\mathcal L$ is,
up to a constant factor, the restriction of a continuous Hermitian metric
on a positive power of an ample line bundle on $\Mn$.
\end{corollary}
\begin{proof}
Grothendieck--Riemann--Roch and Quillen curvature give
\[
 c_1(\mathcal L,h_Q)=-2km^2\,p_*(\delta_v\,\omega)
                   =\frac{rkm^2}{\pi^2}\Om;
 \qquad c_1(\mathcal L)=4m^2c_1(\lambda_D).
\]
These are the surface determinant formulas of
\cite[Theorem~4.6]{BS16}. The current $T$ gives a continuous metric on the
corresponding ample line downstairs, after multiplication by its positive
constant and passage to powers. On $\M$ the ratio of the two metrics is a
flat unitary Hermitian line bundle.

This remaining flat line has finite monodromy. Indeed $H^1(\MG,\QQ)=0$.
The degree-one cohomology sequence for the complement of the boundary
injects $H^1(\M,\QQ)$ into the vector space generated by its divisors,
with next map sending a divisor to its class in $H^2(\MG,\QQ)$.
Those divisors are exceptional, and their classes are linearly independent:
a cohomologically trivial real exceptional divisor is relatively
numerically trivial and vanishes by the negativity lemma, as in
Theorem~\ref{thm:global-current}. Hence $H^1(\M,\QQ)=0$.
The group $H_1(\M,\ZZ)$ is finitely generated, so it is finite.
A further power trivializes the unitary character by a parallel section.
The metrics then agree up to a constant.
\end{proof}

The vanishing used in this proof is specific to the present K3 moduli
space. Equality of curvatures alone would leave a possibly infinite
unitary character on a general surface moduli space. We make no such
continuous Quillen-extension assertion there.

\subsection{Why the Gieseker space does not carry the extended metric}

\begin{corollary}\label{cor:nonextension}
If $\pi$ contracts a curve, $\gL$ has no continuous positive-definite
Riemannian extension to $\MG$.
\end{corollary}
\begin{proof}
Suppose such an extension exists. Since the complex structure is fixed,
continuity makes the extension Hermitian, with a continuous positive
$(1,1)$-form $\widetilde\Omega$. It is closed as a current. To see this,
remove tubular neighbourhoods of the analytic boundary and use cutoffs:
a bounded form paired with the gradient of a cutoff near a real
codimension-two set gives an error tending to zero; higher-codimension
pieces are smaller. Closedness on $\M$ therefore extends across the
boundary.
Locally a continuous closed $(1,1)$-form has a continuous potential $w$;
strict positivity makes $w$ strictly plurisubharmonic. On the complement
of the boundary, $w-\pi^*u$ is pluriharmonic, where $u$ is one of the
potentials of Theorem~\ref{thm:continuous-potential}. Continuity and
removability extend that identity across the boundary. On any contracted
curve, $\pi^*u$ is constant. Thus $w$ restricts locally to a pluriharmonic
function on that curve, contradicting strict plurisubharmonicity.
\end{proof}

This distinguishes the existence of a hyperk\"ahler metric on the smooth
Gieseker moduli space from extension of its particular $L^2$ metric.
Tyurin's extension assertion \cite[Theorem~2]{Tyu92}, made for a fine
component, cannot hold as an extension of this positive-definite metric
when the Uhlenbeck morphism contracts a curve. Example~\ref{ex:eight}
below is fine and has such curves. The obstruction is intrinsic to the
class $\Lambda$, which is zero on every contracted curve.
